% 第 12 章 公司基本面深度卡片
% 每家公司：核心业务、营收结构、库存、现金流、固定资产、利润占比、AI存储敞口

\section{公司基本面全景}

本章为12家核心覆盖标的提供完整的基本面深度卡片，包括核心业务、营收结构、库存变化、现金流质量、资本开支、利润结构和AI存储敞口。\textbf{库存数据是存储涨价周期的核心跟踪指标}——江波龙2025年末存货达116.78亿元（同比+49.1\%），在存储涨价周期中库存增值贡献了显著利润弹性（Q1毛利率55.53\% vs 2025全年19.40\%）。

\begin{riskbox}[基本面研究结论]
\textbf{存储涨价周期的库存红利}：江波龙2025年末存货116.78亿（较年初+49\%），在DRAM/NAND涨价周期中直接转化为毛利增量（Q1毛利率55.53\% vs 2025全年19.40\%）。北方华创存货286亿（设备在手订单饱满），北京君正Q1毛利率43.49\%（存储涨价+车载双轮驱动）。设备厂现金流稳健，模组厂现金流因备货承压但利润弹性最大。澜起科技Q1毛利率69.8\%（受益DDR5高毛利产品占比提升），经营现金流Q1同比+233\%。
\end{riskbox}

\section{核心层公司基本面卡片}

\needspace{46\baselineskip}
\begin{exhibitbox}[表 12-1 \quad 澜起科技 (688008) — DDR5接口芯片全球龙头 / Montage Technology]
\centering
\scriptsize
\renewcommand{\arraystretch}{1.20}
\begin{tabularx}{\exhibitboxwidth}{>{\bfseries}p{2.2cm} >{\raggedright\arraybackslash\sloppy\hspace{0pt}}X}
\toprule
\textbf{维度} & \textbf{详情} \\
\midrule
核心业务 & \textbf{全球领先的集成电路设计企业}，专注云计算和AI基础设施芯片解决方案。三大产品线：(1)内存接口芯片（RCD/DB、MRCD/MDB、CKD、SPD、TS、PMIC），\textbf{全球市占率约36.8\%排名第一}，覆盖DDR2至DDR5全系列；(2)互连类芯片（PCIe Retimer/CXL Retimer），PCIe Retimer全球市占率约10.9\%排名第二；(3)津逮系列平台产品（安全可控CPU）。Fabless模式，晶圆制造主要委托台积电。 \\
产业链地位 & A股AI存储链「一阶Alpha」——每台AI服务器必备RCD，CXL技术布局全球领先（2022年全球首发CXL MXC）。 \\
营收(2025FY) & 总营收\textbf{54.56亿}（YoY+49.9\%）；2026Q1营收14.61亿（YoY+19.5\%）。 \\
营收结构 & 内存接口芯片51.39亿（占94.2\%，YoY+53.4\%）；津逮系列3.08亿（占5.6\%，YoY+10.3\%）；互连类芯片2025Q1营收1.35亿（YoY+155\%，增速最快）。客户以全球三大内存模组厂商（三星、海力士、美光）为主。 \\
库存(2026Q1) & \textbf{8.30亿}（较2025年底有所上升）；存货周转天数从2025年的109天大幅上升至Q1的\textbf{191.5天}。库存增加主因为应对DDR5需求增长主动备货。 \\
现金流 & 2025FY经营CF\textbf{20.22亿}；2026Q1经营CF\textbf{6.27亿}（YoY+232.9\%，回款能力增强）。自由现金流2025年约14.4亿。货币资金146.24亿，资产负债率仅4.17\%，财务安全性极高。 \\
固定资产 & 固定资产\textbf{30.94亿}（2026Q1）；在建工程1.27亿（芯片测试及研发设施）。2025年固定资产投资同比+22.98\%，资本开支加速。 \\
利润结构 & 2025FY归母净利\textbf{22.36亿}（YoY+58.4\%）；2026Q1归母净利\textbf{8.47亿}（YoY+61.3\%）。\textbf{毛利率Q1 69.8\%}（vs 2025FY 62.2\%，受益高毛利DDR5产品占比提升）。内存接口芯片毛利率65.6\%。ROE 2025年18.25\%。 \\
研发投入 & 研发费用率约9.6\%（2025年5.26亿，YoY+168\%，主因股权激励增加）。CXL Switch芯片、GDDR7内存接口芯片在研。 \\
AI存储敞口 & \textbf{极高}——(1)DDR5 RCD全球第一（约36.8\%），已迭代至第五代支持8000MT/s，第二代MRCD/MDB支持12800MT/s；(2)CXL MXC 2025年9月推出CXL3.1版本，支持64GT/s与8000MT/s DDR5；(3)PCIe6.x/CXL3.x Retimer 2025年1月推出支持64GT/s；(4)互连类芯片Q1营收YoY+155\%，反映AI需求爆发。 \\
\bottomrule
\end{tabularx}
\end{exhibitbox}
\sourcenote{财务数据来自公司2025年报（2026-04披露）、2026Q1季报、同花顺F10及AStock研究数据库。业务描述来自公司年报及投资者关系活动记录。}

\needspace{42\baselineskip}
\begin{exhibitbox}[表 12-2 \quad 北方华创 (002371) — 半导体设备平台龙头 / NAURA]
\centering
\scriptsize
\renewcommand{\arraystretch}{1.20}
\begin{tabularx}{\exhibitboxwidth}{>{\bfseries}p{2.2cm} >{\raggedright\arraybackslash\sloppy\hspace{0pt}}X}
\toprule
\textbf{维度} & \textbf{详情} \\
\midrule
核心业务 & \textbf{国内领先的半导体设备平台型企业}，产品覆盖电子工艺装备（93.34\%）和电子元器件（6.55\%）。半导体装备涵盖刻蚀、薄膜沉积（PVD/CVD/ALD）、热处理（炉管）、湿法清洗、离子注入、电镀、涂胶显影（收购芯源微）、混合键合等全系列，是\textbf{国内产品谱系最全的设备厂商}。PVD和立式炉累计出货均突破1000台，刻蚀设备营收超百亿。累计申请专利超11300件。 \\
产业链地位 & A股半导体设备「二阶Beta」——无论AI需求流向哪家原厂，扩产都绕不开北方华创设备。存储厂（长存/长鑫）设备国产化率从10\%$\to$25\%核心受益者。 \\
营收(2025FY) & 总营收\textbf{393.53亿}（YoY+30.85\%）；2026Q1营收103.23亿（YoY+25.80\%）。 \\
营收结构 & 电子工艺装备367.31亿（93.34\%）、电子元器件25.79亿（6.55\%）。集成电路设备营收同比增长超50\%，刻蚀设备和薄膜沉积设备营收均超百亿。区域：中部及东南部245.60亿（62.41\%，增速最快YoY+41.7\%）、东北及华北112.72亿（28.64\%）。 \\
库存(2025年末) & \textbf{286.27亿}（YoY+21.1\%，但占总资产比例从35.61\%降至31.88\%）。存货周转约405天。构成：原材料130.37亿、在产品30.85亿、库存商品124.93亿。2026Q1存货286.03亿，基本持平。高库存反映在手订单饱满（设备交付周期6--12个月）。 \\
现金流 & 2025FY经营CF\textbf{21.33亿}（经营CF/净利=38.6\%，偏低主因行业回款周期长）；2026Q1经营CF 7.48亿（YoY+143\%，大幅改善）。自由现金流约-1.29亿。 \\
固定资产 & 固定资产\textbf{75.68亿}；在建工程4.52亿。2025年购建固定资产支付现金22.62亿（YoY+8.6\%），投资活动现金流出50.54亿（含并购芯源微等20.17亿）。 \\
利润结构 & 2025FY归母净利\textbf{55.22亿}（YoY-1.77\%，扣非YoY-4.20\%）；2026Q1归母净利16.35亿（YoY+3.42\%）。毛利率Q1 40.77\%（vs 2025FY 40.10\%）。电子工艺装备毛利率39.18\%，电子元器件毛利率52.95\%。毛利率同比下降主因新产品验证期零部件迭代升级成本增加。ROE 16.41\%。 \\
AI存储敞口 & \textbf{高}——刻蚀、CVD、PVD、热处理、湿法清洗、电镀等全系列设备已适配3D NAND与HBM制造需求，进入长江存储、长鑫存储等头部存储厂商批量采购清单。TSV设备（深孔刻蚀、ALD、铜种子层PVD、电镀填充）直接受益HBM产能扩张。混合键合设备可应用于HBM及Chiplet领域。2025年集成电路设备营收同比增长超50\%。 \\
\bottomrule
\end{tabularx}
\end{exhibitbox}
\sourcenote{财务数据来自北方华创2025年年报（2026-04-18）、2026Q1季报、东吴证券研报（2026-04-21）及AStock财务数据库。}

\needspace{48\baselineskip}
\begin{exhibitbox}[表 12-3 \quad 江波龙 (301308) — 存储模组涨价弹性最大 / Longsys]
\centering
\scriptsize
\renewcommand{\arraystretch}{1.20}
\begin{tabularx}{\exhibitboxwidth}{>{\bfseries}p{2.2cm} >{\raggedright\arraybackslash\sloppy\hspace{0pt}}X}
\toprule
\textbf{维度} & \textbf{详情} \\
\midrule
核心业务 & \textbf{国内领先的存储模组企业}，全球存储应用市场重要厂商。主要产品：(1)嵌入式存储（UFS/eMMC/ePOP）；(2)固态硬盘SSD（消费级与企业级eSSD）；(3)移动存储（U盘/存储卡/便携式SSD）；(4)内存条（DDR4/DDR5，含企业级RDIMM）。拥有\textbf{Lexar}（全球高端消费存储品牌，19-25年收入CAGR 33\%）、\textbf{Zilia}（巴西/拉美头部存储品牌）两大自主品牌。自研5nm UFS 4.1主控芯片量产，全系列主控累计部署超1.4亿颗。 \\
产业链地位 & \textbf{存储涨价周期最直接受益者}——模组厂在涨价周期中库存增值+ASP提升双击，利润率弹性远大于原厂。2026中报预告验证此逻辑。 \\
营收(2025FY) & 总营收\textbf{227.66亿}（YoY+30.36\%）；2026Q1营收\textbf{99.09亿}（YoY+132.79\%，涨价+放量）。 \\
营收结构 & 嵌入式存储100.12亿（44.0\%，YoY+18.83\%）；固态硬盘55.70亿（24.5\%，YoY+34.31\%）；移动存储48.94亿（21.5\%，YoY+52.56\%）；内存条22.16亿（9.7\%，YoY+45.1\%）。分品牌：Lexar 47.41亿（+34.53\%）；Zilia 29.24亿（+26.49\%）；行业及存储业务17.83亿（+93.30\%）。境外约33.7\%，境内约66.3\%。 \\
库存(2025年末) & \textbf{116.78亿}（较年初+49.1\%，占资产总额70.3\%）。存货周转约187天。\textbf{关键洞察}：存货从78.33亿增至116.78亿，主因为公司预判存储涨价周期主动加大晶圆颗粒备货。2026Q1存储价格加速上涨（NAND/DRAM价格指数2025年9月至2026年4月分别上涨370\%/376\%），低价库存增值成为Q1利润爆发核心驱动。国信预测2026年末存货降至89.68亿（去库存周期开启）。 \\
现金流 & 2025FY经营CF\textbf{-12.01亿}（大幅备货，战略性现金流投入）；自由现金流-20.46亿。\textbf{周期股特征}：涨价初期现金流为负（备货），涨价后期转为强正现金流（去库存）。国信预测2026E经营CF将大幅转正至+130.63亿，自由CF +125.63亿。 \\
固定资产 & 固定资产25.94亿；其他长期资产（含在建工程）合计30.87亿（YoY+27.9\%）。2025FY CAPEX 8.45亿。国信预测2026E CAPEX 5.0亿、2027E 3.0亿，呈逐年下降趋势，产能扩建周期接近尾声。 \\
利润结构 & 2025FY净利\textbf{14.23亿}（YoY+185.41\%）；2026Q1净利\textbf{38.62亿}（YoY+2644\%，\textbf{单季已超去年全年}）。\textbf{毛利率Q1 55.53\%}（vs 2025FY 19.40\%，+36pp！），净利率40.16\%，均创历史新高。 \\
\midrule
\rowcolor{riskgreen!8}
\textbf{中报预告} & \textbf{2026H1净利92--110亿（YoY+62204\%~+74394\%）}，Q2单季53.4--71.4亿（QoQ+38\%~+84\%）。存储涨价周期盈利弹性充分验证。 \\
\midrule
AI存储敞口 & \textbf{极高}——企业级存储2025年收入17.83亿（+93.3\%），占总营收7.8\%。已构建"eSSD+RDIMM+SOCAMM+MRDIMM+CXL2.0"全栈企业级产品体系，eSSD与RDIMM完成鲲鹏、海光、飞腾、AMD等主流平台兼容认证，切入通信、金融、互联网头部客户。国信预测2026年企业级SSD收入有望达80亿。HBM模组预计2026年量产，初期贡献<3\%。端侧AI存储：自研5nm UFS 4.1主控量产，SPU/iSA/HLC技术形成差异化壁垒。 \\
\bottomrule
\end{tabularx}
\end{exhibitbox}
\sourcenote{财务数据来自江波龙2025年度报告（深交所2026-04-26披露）、2026年中报业绩预告（2026-07-03公告）、国信证券《江波龙1Q26归母净利润同比增长2644.05\%》（2026-05-07）及爱建证券年报点评。}

\needspace{44\baselineskip}
\begin{exhibitbox}[表 12-4 \quad 中微公司 (688012) — 刻蚀设备龙头 / AMEC]
\centering
\scriptsize
\renewcommand{\arraystretch}{1.20}
\begin{tabularx}{\exhibitboxwidth}{>{\bfseries}p{2.2cm} >{\raggedright\arraybackslash\sloppy\hspace{0pt}}X}
\toprule
\textbf{维度} & \textbf{详情} \\
\midrule
核心业务 & \textbf{国内领先的微观加工高端设备企业}。核心产品：(1)刻蚀设备：CCP刻蚀（Primo D-RIE等9款，累计出货超5000反应台）和ICP刻蚀（Primo AD-RIE等，累计装机1800反应台），技术达国际先进水平；(2)MOCVD设备：PRISMO D-BLUE等4款，中国大陆LED MOCVD市场占有率第一；(3)薄膜沉积设备：LPCVD、ALD等6款；(4)环保设备。\textbf{国内唯一能与国际巨头竞争的刻蚀设备厂商}。 \\
产业链地位 & 全球刻蚀设备Top5（唯一亚洲厂商）；CCP刻蚀国内份额第一，ICP刻蚀快速突破。 \\
营收(2025FY) & 总营收\textbf{123.85亿}（YoY+36.62\%）；2026Q1营收29.15亿（YoY+34.13\%）。 \\
营收结构 & 刻蚀设备约81亿（占约65\%）；MOCVD约18亿（占约15\%）；薄膜沉积约7亿（占约5.5\%，前三季度YoY+1333\%）；其他约18亿（占约14.5\%）。估算境内约85-90\%、境外约10-15\%。 \\
库存(2025年末) & \textbf{20.64亿}（随业务规模扩大持续增长）。存货周转天数从年初513天改善至\textbf{339天}（2025年），2026Q1约370天，周转效率提升明显。 \\
现金流 & 2025FY经营CF约\textbf{22.8亿}（经营CF/净利约108\%，现金流质量较好）；2026Q1经营CF -1.59亿（季节性波动，主因支付供应商货款及职工薪酬增加）。自由现金流约15-18亿。 \\
固定资产 & 固定资产约4.78亿（2025年末）；在建工程约1.26亿。持续推进产能扩张，2025年研发费用约37.36亿（YoY+52.32\%）。 \\
利润结构 & 2025FY归母净利\textbf{21.11亿}（YoY+30.69\%，扣非YoY+11.64\%）；2026Q1归母净利9.30亿（YoY+197.20\%，主因投资收益等非经常性损益增加，扣非YoY+60.09\%）。毛利率Q1 39.89\%（vs 2025FY 39.17\%）。刻蚀设备毛利率约40-42\%，为核心毛利来源（约65-70\%）。ROE 9.97\%。 \\
AI存储敞口 & \textbf{高}——CCP刻蚀设备在存储芯片领域已有批量应用，客户涵盖长江存储等国内主要存储芯片厂商（估算存储相关刻蚀需求占刻蚀设备收入约20-30\%）。3D NAND叠层刻蚀（长存232L/290L）、DRAM高深宽比刻蚀（长鑫17nm/15nm）、HBM TSV刻蚀均为核心设备需求。ICP刻蚀可用于先进封装领域，间接服务AI芯片封装需求。 \\
\bottomrule
\end{tabularx}
\end{exhibitbox}
\sourcenote{财务数据来自中微公司2025年业绩快报、2026Q1季报及公司官网公开信息。}

\needspace{44\baselineskip}
\begin{exhibitbox}[表 12-5 \quad 拓荆科技 (688072) — ALD/CVD设备龙头 / Piotech]
\centering
\scriptsize
\renewcommand{\arraystretch}{1.20}
\begin{tabularx}{\exhibitboxwidth}{>{\bfseries}p{2.2cm} >{\raggedright\arraybackslash\sloppy\hspace{0pt}}X}
\toprule
\textbf{维度} & \textbf{详情} \\
\midrule
核心业务 & \textbf{国内领先的半导体薄膜沉积设备厂商}。核心产品覆盖PECVD、ALD、SACVD、HDPCVD、FlowableCVD等系列薄膜沉积设备，以及三维封装系列设备（晶圆直接键合、临时键合、减薄、清洗等）。\textbf{国产CVD/ALD设备的核心供应商}，产品广泛应用于逻辑芯片、存储芯片（DRAM/NAND/3D NAND）、功率器件等领域。累计申请专利2140项（含PCT国际专利707项）。 \\
产业链地位 & 长存/长鑫薄膜沉积设备核心供应商；HBM制造中ALD是关键制程（高介电常数材料沉积）。 \\
营收(2025FY) & 总营收\textbf{65.19亿}（YoY+58.87\%，增速最快设备厂）；2026Q1营收11.12亿（YoY+56.97\%）。 \\
营收结构 & PECVD约65-70\%（核心主力产品）；ALD约15-20\%（快速增长，2025Q2新增订单超2024全年）；SACVD及其他约10-15\%。中国大陆约92-95\%（核心客户中芯国际、长江存储、长鑫存储）；海外约5-8\%。 \\
库存(2025年末) & \textbf{约9-12亿}（2025年中报8.80亿，YoY+32.85\%）。存货周转天数较长（2025年约639天，2026Q1约1112天），反映半导体设备行业长周期生产验证特性。2026Q1存货跌价准备仅580万元（较年初-56.88\%），\textbf{库存质量健康}。 \\
现金流 & 2025FY经营CF\textbf{35.9亿}（经营CF/净利约3.87，\textbf{现金流质量优秀}，主因设备预收款及回款节奏改善）；2026Q1经营CF -5.20亿（季节性）。自由现金流约16-21亿。 \\
固定资产 & 固定资产16.91亿（2026Q1，较年初+89.34\%，因联合验收装备类项目转入）；在建工程4.13亿（2026Q1，较年初-43.32\%）。沈阳总部产能扩建推进中。 \\
利润结构 & 2025FY归母净利\textbf{9.27亿}（YoY+34.67\%，扣非YoY+103.05\%）；2026Q1归母净利5.71亿（YoY+488.29\%，主因高毛利设备交付结构优化及投资收益）。\textbf{毛利率Q1 41.69\%}（vs 2025FY 34.95\%，+6.7pp，主因高毛利ALD设备占比提升）。ROE 15.77\%。 \\
\midrule
\rowcolor{riskamber!8}
\textbf{重大事项} & \textbf{2026-06-29起停牌}：筹划发行股份及支付现金收购\textbf{无锡尚积半导体}控股权，旨在补上PVD与刻蚀拼图，平台化布局加速。 \\
\midrule
AI存储敞口 & \textbf{高}——(1)ALD设备在新型存储器（含HBM相关制程）及AI相关领域应用持续突破，2025Q2 ALD新增订单超2024全年；(2)长江存储为核心客户之一，ALD设备来自长存收入占比从2024年约12\%提升至2025年约18\%；(3)SACVD设备已通过中芯国际5nm制程验证；(4)HBM封装TSV薄膜沉积设备占国产设备投资约8\%（拓荆与北方华创共同覆盖）。 \\
\bottomrule
\end{tabularx}
\end{exhibitbox}
\sourcenote{财务数据来自拓荆科技2025年报（上交所SSE, 2026-04-26）及AIDC供应链研究数据。停牌信息来自公司公告（2026-06-29）。}

\section{卫星层公司基本面卡片}

\needspace{46\baselineskip}
\begin{exhibitbox}[表 12-6 \quad 北京君正 (300223) — 车载存储+利基DRAM / ISSI]
\centering
\scriptsize
\renewcommand{\arraystretch}{1.20}
\begin{tabularx}{\exhibitboxwidth}{>{\bfseries}p{2.2cm} >{\raggedright\arraybackslash\sloppy\hspace{0pt}}X}
\toprule
\textbf{维度} & \textbf{详情} \\
\midrule
核心业务 & \textbf{Fabless集成电路设计公司}，坚持"存储+计算+模拟"产品战略。核心产品线：(1)存储芯片：SRAM、DRAM、Nor Flash，\textbf{车规级存储全球领先，利基DRAM全球第六/中国大陆第一，车规SRAM全球第一}；(2)计算芯片：基于MIPS/RISC-V架构，面向智能安防、AIoT、泛视觉市场；(3)模拟与互联芯片：LED驱动、车规互联芯片。ISSI（原美国上市公司，2015年被君正收购）是全球利基DRAM龙头。 \\
产业链地位 & ISSI在汽车/工业/医疗领域份额领先。AI驱动存储超级周期：三星/美光/海力士将产能转向HBM，传统DRAM产能被挤压引发涨价潮，公司作为利基DRAM龙头直接受益。 \\
营收(2025FY) & 总营收\textbf{47.41亿}（YoY+12.54\%）；2026Q1营收15.60亿（YoY+47.12\%，存储涨价+车载放量）。 \\
营收结构 & 存储芯片29.11亿（61.40\%）、计算芯片12.93亿（27.28\%）、模拟与互联5.06亿（10.67\%）。区域：境外39.19亿（82.67\%）、境内8.13亿（17.15\%）。 \\
库存(2025年末) & \textbf{29.75亿}（账面余额31.79亿，跌价准备2.04亿，较年初+11.3\%）。库存量95,855万颗（YoY+47.58\%），主因预判存储超级周期加大备货。存货周转约330天，处于Fabless设计公司合理水平。 \\
现金流 & 2025FY经营CF\textbf{6.23亿}（YoY+71.30\%，OCF/营收比13.13\%，高于净利率7.92\%）；2026Q1经营CF 2.17亿（YoY+23.72\%）。自由现金流约5.5亿。\textbf{无短期和长期借款，资产负债率仅7.74\%，财务极为稳健}。 \\
固定资产 & 固定资产4.62亿（Fabless模式，主要为研发设备）；在建工程0.10亿。2025年CAPEX约0.7亿，无晶圆厂建设计划。 \\
利润结构 & 2025FY净利\textbf{3.76亿}（YoY+2.74\%）；2026Q1净利\textbf{3.19亿}（YoY+331.61\%，\textbf{单季接近去年全年}）。\textbf{毛利率Q1 43.49\%}（vs 2025FY 34.09\%，+9.4pp）。分业务：存储芯片毛利率38.89\%（Q1，vs 2025年31.95\%）；计算芯片51.90\%（Q1）；模拟互联50.88\%（Q1）。存储芯片毛利贡献约58\%。ROE 3.06\%。 \\
\midrule
\rowcolor{riskgreen!8}
\textbf{调研确认} & \textbf{2026-07-03机构调研}：存储芯片持续在涨价，客户基本接受；预计各季度毛利率环比有所增长；DRAM代工产能改善窗口预计在2027年下半年；2027年供需仍紧张。 \\
\midrule
AI存储敞口 & \textbf{中高}——(1)AI驱动存储超级周期：HBM产能挤压传统DRAM产能，利基DRAM涨价潮中公司直接受益；(2)3D DRAM研发启动：面向AI应用的3D DRAM中算力SoC和DRAM接口单元base die研发，通过混合键合技术叠加高带宽、低功耗、端侧AI能力；(3)AI计算芯片：基于RISC-V，NPU即将迭代至4.0（512Tops）；(4)AI MCU：首款具备AI算力的高性能MCU已完成研发与样品投片。存储芯片Q1收入10.18亿（YoY+53.6\%）。 \\
\bottomrule
\end{tabularx}
\end{exhibitbox}
\sourcenote{财务数据来自公司2025年年度报告、2026年一季度报告、中邮证券研究报告（2025-12-25）、华安证券研究报告（2025-10-30）及AIDC供应链研究数据。调研信息来自2026-07-03机构调研纪要。}

\needspace{42\baselineskip}
\begin{exhibitbox}[表 12-7 \quad 沪硅产业 (688126) — 硅片材料龙头 / Zing Semi]
\centering
\scriptsize
\renewcommand{\arraystretch}{1.20}
\begin{tabularx}{\exhibitboxwidth}{>{\bfseries}p{2.2cm} >{\raggedright\arraybackslash\sloppy\hspace{0pt}}X}
\toprule
\textbf{维度} & \textbf{详情} \\
\midrule
核心业务 & \textbf{国内半导体硅片龙头企业}。核心产品：(1)300mm半导体硅片（子公司上海新昇、晋科硅材料）；(2)200mm及以下尺寸硅片（新傲科技、芬兰Okmetic）；(3)SOI硅片（新傲科技、新傲芯翼）；(4)单晶压电薄膜衬底（新硅聚合）。上海+太原300mm合计产能\textbf{85万片/月}，承担国家02专项，实现300mm硅片逻辑与存储工艺全覆盖、国内主要客户全覆盖。 \\
产业链地位 & 国内唯一实现300mm硅片规模化销售的厂商；上海新昇是大硅片国产化主力。 \\
营收(2025FY) & 总营收\textbf{37.16亿}（YoY+9.69\%）；2026Q1营收10.84亿（YoY+35.22\%，加速增长）。 \\
营收结构 & 300mm硅片24.39亿（66.3\%）；200mm及以下11.25亿（30.6\%）；受托加工1.15亿（3.1\%）。区域：中国境内26.73亿（72.7\%），其他国家和地区10.06亿（27.3\%）。 \\
库存(2025年末) & \textbf{19.66亿}（较年初+27.5\%，主因300mm硅片产能扩张导致库存量增加，300mm硅片库存量182.23万片YoY+221.65\%）。存货跌价准备3.42亿（占存货14.82\%，反映行业价格下行压力）。2026Q1约20.49亿。存货周转约160天。 \\
现金流 & 2025FY经营CF\textbf{-5.59亿}（较2024年-7.88亿减亏）；自由现金流约-54.68亿。\textbf{重资产+爬坡期}特征：经营现金流尚不能覆盖资本开支，依赖外部融资（2025年筹资活动现金净流入64.72亿）。 \\
固定资产 & \textbf{固定资产140.69亿}（重资产！）；在建工程53.28亿。2025年购建固定资产支付现金49.09亿。已签约未列示的资本性支出承诺33.41亿。太原项目持续建设。 \\
利润结构 & 2025FY净亏损\textbf{15.08亿}（YoY-55.33\%，亏损扩大）；2026Q1亏损4.83亿（YoY-131.67\%）。毛利率Q1 \textbf{-11.85\%}（vs 2025FY -17.06\%）。分业务：300mm硅片毛利率-14.99\%，200mm及以下-23.37\%。整体亏损主因产能爬坡阶段固定成本高、规模效应未充分释放，以及商誉减值（约4亿）和存货跌价损失增加。\textbf{PE估值屏蔽}，采用PS/PB估值。ROE -12.68\%。 \\
\midrule
\rowcolor{riskamber!8}
\textbf{扩产计划} & 上海+太原300mm合计产能已达85万片/月，太原基地产能释放将进一步支撑存储芯片客户需求。重点研发面向AI/数据中心的300mm高规格硅片。 \\
\midrule
AI存储敞口 & \textbf{中}——300mm硅片已通过长江存储（3D NAND）、长鑫存储（DRAM）认证并持续批量供货，\textbf{存储类抛光片占产品比例约60-65\%}。300mm硅片业务实现存储工艺产品全覆盖、国内主要存储客户全覆盖。行业层面HBM抛光片需求强劲，公司300mm硅片可应用于DRAM、3D NAND等存储芯片制造。正重点研发面向AI/数据中心应用的300mm高规格硅片。 \\
\bottomrule
\end{tabularx}
\end{exhibitbox}
\sourcenote{财务数据来自沪硅产业2025年年度报告（上交所披露）、同花顺F10（basic.10jqka.com.cn/688126）及公司投资者互动平台公开信息。}

\needspace{42\baselineskip}
\begin{exhibitbox}[表 12-8 \quad 深科技 (000021) — 封测+存储模组 / Kaifa]
\centering
\scriptsize
\renewcommand{\arraystretch}{1.20}
\begin{tabularx}{\exhibitboxwidth}{>{\bfseries}p{2.2cm} >{\raggedright\arraybackslash\sloppy\hspace{0pt}}X}
\toprule
\textbf{维度} & \textbf{详情} \\
\midrule
核心业务 & \textbf{全球领先的专业电子制造企业}，构建存储半导体、高端制造、计量智能终端三大主业格局。存储半导体：\textbf{国内高端存储芯片封测龙头}，从事DRAM、NAND Flash及嵌入式存储芯片的封装与测试（DDR4/DDR5、LPDDR4/LPDDR5、UFS、uMCP），深圳沛顿+合肥沛顿双基地运营。高端制造：深耕医疗健康、汽车电子、消费电子等。累计授权专利超2800项。 \\
产业链地位 & 国内存储封测重要玩家；子公司深圳沛顿为广东省制造业单项冠军，合肥沛顿为国家专精特新小巨人。 \\
营收(2025FY) & 总营收\textbf{157.47亿}（YoY+6.21\%）；2026Q1营收37.24亿（YoY+10.67\%）。 \\
营收结构 & 高端制造85.35亿（54.20\%）、存储半导体40.91亿（25.98\%，YoY+16.16\%）、计量智能终端30.21亿（19.18\%）。区域：亚太（除中国）75.59亿（48.00\%）、中国（含香港）58.20亿（36.96\%）。 \\
库存(2026Q1) & \textbf{28.67亿}（较2025年末+18.1\%，较2024年末+9.0\%）。存货周转约65-72天。存货增长主要与存储封测业务订单扩大、原材料备货增加有关。 \\
现金流 & 2025FY经营CF\textbf{17.65亿}（YoY-27.3\%，主因购买商品支付现金增加）；2026Q1经营CF 1.14亿（YoY-82.9\%）。经营CF/净利=1.22，现金流质量良好。自由现金流约5.45亿。 \\
固定资产 & 固定资产64.08亿；在建工程0.34亿。2025年CAPEX 12.20亿（主要用于存储封测产能扩充及设备采购）。深圳、合肥双基地持续扩产。 \\
利润结构 & 2025FY归母净利\textbf{11.36亿}（YoY+22.07\%）；2026Q1归母净利2.42亿（YoY+35.35\%）。毛利率Q1 17.07\%（vs 2025FY 18.32\%）。分业务：存储半导体毛利率20.40\%（收入占比25.98\%）；计量智能终端39.26\%；高端制造9.71\%。ROE 9.09\%。 \\
AI存储敞口 & \textbf{中低}——存储半导体业务为AI存储核心暴露环节，2025年营收40.91亿（占25.98\%）。产品覆盖DRAM(DDR4/DDR5)、LPDDR4/LPDDR5、NAND Flash(UFS3.1/4.1)、嵌入式存储(uMCP)等。技术方面具备NAND Flash 32D超高堆叠封装和GDDR多芯片倒装封装研发能力，LPDDR5 PoPt超薄叠层封装方案已通过客户认证。公司明确提出"加强在AI驱动下的先进存储封装布局"战略，成立先进封装研究院。HBM作为AI存储核心增量引擎，DDR5在服务器端渗透率快速提升，公司正积极把握机遇。 \\
\bottomrule
\end{tabularx}
\end{exhibitbox}
\sourcenote{财务数据来自深科技2025年年度报告（巨潮资讯网2026-04-29）、2026年Q1报告及AStock研究数据。}

\needspace{42\baselineskip}
\begin{exhibitbox}[表 12-9 \quad 长电科技 (600584) — 先进封测龙头 / JCET]
\centering
\scriptsize
\renewcommand{\arraystretch}{1.20}
\begin{tabularx}{\exhibitboxwidth}{>{\bfseries}p{2.2cm} >{\raggedright\arraybackslash\sloppy\hspace{0pt}}X}
\toprule
\textbf{维度} & \textbf{详情} \\
\midrule
核心业务 & \textbf{全球第三大半导体封测厂商（OSAT Top3）}，中国封测行业龙头。主要产品：(1)先进封装：SiP系统级封装、Fan-out扇出封装、2.5D/3D封装（XDFOI/Chiplet）、Bumping凸块、WLCSP晶圆级芯片尺寸封装；(2)传统封装：QFN、BGA、SOP等系列；(3)测试服务：晶圆测试、成品测试等。全球OSAT第三（仅次于ASE、Amkor），中国大陆第一。 \\
产业链地位 & 全球OSAT Top3；2.5D/3D封装是HBM必经之路，但TSMC垄断CoWoS约80\%份额，长电产能和客户认证仍在建设中。 \\
营收(2025FY) & 总营收\textbf{388.71亿}（YoY+8.09\%）；2026Q1营收91.71亿（YoY-1.76\%，季节性因素+产品结构调整）。 \\
营收结构 & 先进封装约45\%（SiP/Fan-out/2.5D/3D/Bumping/WLCSP）；传统封装约40\%（QFN/BGA/SOP等）；测试服务约15\%。区域：境外约65\%（核心客户高通、苹果、AMD等国际大厂）；境内约35\%。 \\
库存(2026Q1) & \textbf{44.08亿}。存货周转约47-65天（封测行业因代工模式存货周转快于设计和设备行业）。库存/营收比约0.48x，处于合理区间。未见明显积压或去库迹象。 \\
现金流 & 2025FY经营CF\textbf{46.52亿}（YoY-20.26\%，主因新建工厂产能爬坡期投入加大、原材料成本上升及财务费用增加）；2026Q1经营CF 17.79亿（YoY+55.44\%，显著改善）。2026年CAPEX计划99.8亿，自由现金流仍承压。 \\
固定资产 & \textbf{固定资产约180亿}（重资产模式）；在建工程约35亿。2026年6月宣布拟投资\textbf{78亿}在上海临港建设高端先进封测工厂，聚焦AI芯片先进封装。张江研发大楼已投入使用。 \\
利润结构 & 2025FY归母净利\textbf{15.65亿}（YoY-2.75\%）；2026Q1归母净利2.90亿（YoY+42.74\%，扣非YoY+37.03\%）。毛利率Q1 14.55\%（vs 2025FY 14.15\%，持续改善主因产品结构优化）。先进封装（约45\%营收）毛利率高于传统封装，是利润主要贡献来源。ROE 5.46\%。 \\
\midrule
\rowcolor{riskamber!8}
\textbf{扩产计划} & 拟投资\textbf{78亿}在上海临港建设高端先进封测工厂，聚焦AI芯片先进封装。2026年固定资产投资计划99.8亿元。 \\
\midrule
AI存储敞口 & \textbf{中}——(1)HBM必须经过2.5D/3D先进封装，公司具备XDFOI/Chiplet技术能力，但当前HBM封装收入占比仍低（<5\%）；(2)AI芯片（GPU/ASIC）封装是主要增量来源，核心客户包括高通、AMD等，临港78亿项目聚焦AI芯片先进封装；(3)当前AI相关收入实际贡献<3\%，需2027年H2后方能体现CoWoS类产能贡献。 \\
\bottomrule
\end{tabularx}
\end{exhibitbox}
\sourcenote{财务数据来自长电科技2025年报+2026Q1季报、华源证券研报（2026-05-07）、华鑫证券研报（2026-05-12）及公司临港项目公告（2026-06-25）。}

\needspace{40\baselineskip}
\begin{exhibitbox}[表 12-10 \quad 通富微电 (002156) — AMD封测核心 / Tongfu Microelectronics]
\centering
\scriptsize
\renewcommand{\arraystretch}{1.20}
\begin{tabularx}{\exhibitboxwidth}{>{\bfseries}p{2.2cm} >{\raggedright\arraybackslash\sloppy\hspace{0pt}}X}
\toprule
\textbf{维度} & \textbf{详情} \\
\midrule
核心业务 & \textbf{国内领先的集成电路封装测试企业，AMD封测独家战略合作伙伴}。主要产品：处理器封测（CPU/GPU/APU，AMD为核心客户，收入占比约50\%）、存储芯片封测（DRAM/NAND Flash）、功率器件封测、显示驱动芯片封测。先进封装技术：FCBGA、SiP、Chiplet 2D+、3nm封装已通过AMD认证，CPO在研。全球封测行业排名前列，国内封测三强之一。累计申请专利超1779项。 \\
产业链地位 & AMD MI300/MI400系列AI GPU封测核心供应商；AMD深度联合研发是最大竞争优势。 \\
营收(2025FY) & 总营收\textbf{279.21亿}（YoY+16.92\%）；2026Q1营收74.82亿（YoY+22.80\%）。 \\
营收结构 & 处理器封测约55\%（AMD为主）；存储封测约20\%；功率器件/显示驱动/其他约25\%。区域：境外约55\%（主要来自AMD等海外客户）；境内约45\%。 \\
库存(2026Q1) & \textbf{48.11亿}。存货周转约63天。存货规模与营收增长基本匹配，周转效率保持稳定。 \\
现金流 & 2025FY经营CF\textbf{69.66亿}（YoY+79.66\%，\textbf{经营CF/净利约5.7倍，现金流质量显著优于利润表表现}）；2026Q1经营CF 9.42亿。自由现金流约20-25亿。 \\
固定资产 & 固定资产约120亿；在建工程约15亿。2025年计划各基地含研发在内总投入60亿；2026年计划总投入91亿。 \\
利润结构 & 2025FY归母净利\textbf{12.19亿}（YoY+79.86\%）；2026Q1归母净利3.29亿（YoY+224.55\%）。毛利率Q1 13.32\%（vs 2025FY 14.59\%）。处理器封测（AMD）贡献主要毛利（毛利率约14-16\%）；存储封测毛利随存储周期波动。ROE 8.08\%。 \\
\midrule
\rowcolor{riskgreen!8}
\textbf{定增获批} & \textbf{2026-07-03}：42.2亿定增获证监会同意注册批复，用于存储芯片封测等项目。国家大基金持股已降至5\%以下。 \\
\midrule
AI存储敞口 & \textbf{中低（约5-12\%）}——核心增量来自AMD AI GPU封测：作为AMD独家封测战略伙伴，AMD MI300/MI400系列AI芯片封测订单直接拉动FCBGA等高端封装业务。存储芯片封测间接受益于AI服务器存储升级（DDR5容量倍增），但HBM封装非核心业务。2026年7月42.2亿定增获批，用于存储芯片封测等项目，强化存储封测产能布局。潜在客户包括海光信息、寒武纪、华为昇腾等。 \\
\bottomrule
\end{tabularx}
\end{exhibitbox}
\sourcenote{财务数据来自通富微电2025年年报（2026-04-24）、2026Q1季报及2026-07-03定增获批公告。}

\section{主题层公司基本面卡片}

\needspace{46\baselineskip}
\begin{exhibitbox}[表 12-11 \quad 兆易创新 (603986) — NOR Flash+利基DRAM / GigaDevice]
\centering
\scriptsize
\renewcommand{\arraystretch}{1.20}
\begin{tabularx}{\exhibitboxwidth}{>{\bfseries}p{2.2cm} >{\raggedright\arraybackslash\sloppy\hspace{0pt}}X}
\toprule
\textbf{维度} & \textbf{详情} \\
\midrule
核心业务 & \textbf{全球领先的Fabless芯片供应商}，形成"感存算控连"多元产品矩阵。核心产品：(1)NOR Flash：\textbf{全球第二大供应商（市占率约18.5\%）}，覆盖512Kb至2Gb全容量，45nm节点SPI NOR Flash大规模量产；(2)利基DRAM：DDR3/DDR4/LPDDR全产品线布局，DDR4 8Gb快速抢占份额，自研LPDDR4系列研发中；(3)MCU：\textbf{国内32位通用MCU龙头}，69大系列超700款产品，覆盖ARM和RISC-V内核；(4)传感器：触控芯片、指纹识别芯片、气压传感器；(5)模拟芯片：DC-DC、LDO、电机驱动等（2024年底收购苏州赛芯补强）。 \\
产业链地位 & NOR Flash全球Top3（仅次于旺宏、华邦电）；利基DRAM国内份额第一；MCU国内32位MCU龙头。 \\
营收(2025FY) & 总营收\textbf{92.03亿}（YoY+25.12\%）；2026Q1营收\textbf{41.88亿}（YoY+119.38\%，存储涨价+放量）。 \\
营收结构 & 存储芯片65.66亿（71.3\%，含NOR Flash约55-60\%、利基DRAM约30-35\%、SLC NAND约5-10\%，毛利率42.84\%）；MCU 19.10亿（20.8\%，毛利率35.82\%）；传感器3.89亿（4.2\%）；模拟产品3.33亿（3.6\%，毛利率36.96\%）。区域：中国大陆27.87亿（30.3\%）；中国香港43.20亿（46.9\%，含转口贸易）；其他国家及地区20.95亿（22.8\%，毛利率53.75\%）。 \\
库存(2026Q1) & \textbf{34.01亿}（较年初+10.9\%，YoY+38.34\%）。存货周转2025年约177天。库存增加主因存储芯片备货增加（库存量+19.35\%）和MCU库存增加（库存量+65.75\%）。整体库存水平随行业景气度回升而增加，处于合理备货区间。 \\
现金流 & 2025FY经营CF\textbf{21.29亿}（YoY+4.74\%，经营CF/净利=1.29，现金流质量良好）；2026Q1经营CF\textbf{17.83亿}（YoY+430.91\%，\textbf{现金流爆发}）。2025年自由现金流约10.80亿。2025FY CAPEX 10.49亿（YoY+110\%）。 \\
固定资产 & 固定资产14.06亿（2026Q1，Fabless模式，以研发设备和办公资产为主）；在建工程0.16亿。 \\
利润结构 & 2025FY归母净利\textbf{16.48亿}（YoY+49.47\%）；2026Q1归母净利\textbf{14.61亿}（YoY+522.79\%，\textbf{单季接近去年全年}）。\textbf{毛利率Q1 57.08\%}（vs 2025FY 40.22\%，+16.9pp！存储涨价驱动）。分业务：存储芯片毛利率42.84\%（+2.57pp）；MCU 35.82\%（-0.92pp）；传感器19.57\%（+3.11pp）；模拟36.96\%（+26.43pp）。ROE 2025年9.30\%。 \\
\midrule
\rowcolor{riskamber!8}
\textbf{风险提示} & \textbf{2026-06-29发布风险提示}：股价短期涨幅较大，提醒投资者注意交易风险。7月2日跌停。 \\
\midrule
AI存储敞口 & \textbf{中}——四大方向：(1)NOR Flash：端侧AI核心代码存储载体，全球市占率约18.5\%排名第二，AI眼镜、AIPC、智能穿戴等终端AI化驱动需求扩容；(2)利基DRAM：受益AI服务器HBM/DDR5需求爆发，海外大厂退出利基市场，DDR4 8Gb快速抢占份额，自研LPDDR4预计2026年量产；(3)定制化存储：控股子公司青耘科技布局定制化存储，2026年有望在汽车座舱、AIPC、机器人等领域突破；(4)AI MCU：GD32H7系列MCU支持AI算法，无线MCU支持AI大语言模型方案。整体AI相关收入目前以NOR Flash和利基DRAM为主。 \\
\bottomrule
\end{tabularx}
\end{exhibitbox}
\sourcenote{财务数据来自兆易创新2025年年报（上交所2026-04）、2026Q1季报及太平洋证券深度报告（2025-12）。风险提示来自公司公告（2026-06-29）。}

\needspace{40\baselineskip}
\begin{exhibitbox}[表 12-12 \quad 南大光电 (300346) — 半导体材料/前驱体 / Nanda Optoelectronic]
\centering
\scriptsize
\renewcommand{\arraystretch}{1.20}
\begin{tabularx}{\exhibitboxwidth}{>{\bfseries}p{2.2cm} >{\raggedright\arraybackslash\sloppy\hspace{0pt}}X}
\toprule
\textbf{维度} & \textbf{详情} \\
\midrule
核心业务 & \textbf{国内MO源（金属有机源）龙头企业}。主要产品：(1)MO源——LED和半导体制造关键前驱体材料，\textbf{国内市占率超60\%}；(2)ALD/CVD前驱体——应用于DRAM/HBM等存储芯片和逻辑芯片制造，\textbf{IC客户已成为第一大营收来源（占比超30\%）}；(3)电子特气——含氟电子特气系列；(4)ArF浸没式光刻胶——研发中，预计2027年客户验证。技术壁垒在于高纯金属有机化合物合成、纯化和分析检测技术，属于半导体材料卡脖子环节。 \\
产业链地位 & MO源国内市场份额约60\%；ALD前驱体是HBM/先进制程制造的关键材料。 \\
营收(2025FY) & 总营收\textbf{25.85亿}（YoY+9.93\%）；2026Q1营收6.62亿（YoY+5.45\%，增速放缓，前驱体高增长被MO源周期性波动部分抵消）。 \\
营收结构 & MO源约45\%；前驱体（含ALD/CVD用）约25\%；电子特气约20\%；其他约10\%。区域：境内约80\%；境外约20\%。 \\
库存(2026Q1) & \textbf{8.11亿}。存货周转天数从2025年的约180天上升至Q1的约194天，库存水平有所增加，与前驱体业务备货和产能扩张有关。 \\
现金流 & 2025FY经营CF\textbf{7.33亿}（YoY+16.76\%，经营CF/净利=2.29x，\textbf{盈利质量高}）；2026Q1经营CF 1.00亿（YoY-39.24\%，季节性营运资金波动）。自由现金流约5.63亿。 \\
固定资产 & 固定资产约15亿；在建工程约5亿（宁波/南通基地）。2025年CAPEX 1.70亿，主要用于宁波/南通生产基地扩产和前驱体产能建设。 \\
利润结构 & 2025FY净利\textbf{3.20亿}（YoY+18.00\%，扣非YoY+31.49\%）；2026Q1净利1.24亿（YoY+29.97\%，扣非YoY+38.61\%）。毛利率Q1 41.86\%（vs 2025FY 39.62\%，同比提升1.0pp）。前驱体业务2025H1毛利同比增长超50\%（主要增长引擎）。ROE 9.19\%。 \\
AI存储敞口 & \textbf{中低}——前驱体材料是HBM/DRAM晶圆制造的必需材料，公司作为国内MO源龙头和ALD/CVD前驱体核心供应商，间接受益于AI存储扩产。HBM前驱体国产份额约2\%（保守2026E），HBM前驱体业务约占营收5\%。IC客户已成为第一大营收来源（占比超30\%），前驱体业务2025H1毛利同比增长超50\%。整体AI存储相关收入占比仍低，先进制程前驱体份额有待提升。ArF浸没式光刻胶研发中，预计2027年客户验证。 \\
\bottomrule
\end{tabularx}
\end{exhibitbox}
\sourcenote{财务数据来自南大光电2025年年报、2026Q1季报（巨潮资讯网）及AStock研究数据库。}

\section{跨公司关键指标对比}

\begin{exhibitbox}[表 12-13 \quad 12家核心标的关键财务指标横向对比 / Cross-Company Financial Comparison]
\centering
\tiny
\renewcommand{\arraystretch}{1.18}
\setlength{\tabcolsep}{2pt}
\begin{tabularx}{\exhibitboxwidth}{>{\bfseries\raggedright\arraybackslash}p{1.3cm} >{\raggedleft\arraybackslash}p{0.9cm} >{\raggedleft\arraybackslash}p{0.8cm} >{\raggedleft\arraybackslash}p{0.8cm} >{\raggedleft\arraybackslash}p{0.8cm} >{\raggedleft\arraybackslash}p{0.8cm} >{\raggedleft\arraybackslash}p{0.8cm} >{\raggedleft\arraybackslash}p{0.8cm} >{\raggedleft\arraybackslash}p{0.8cm} >{\raggedright\arraybackslash\sloppy\hspace{0pt}}X}
\toprule
\textbf{公司} & \textbf{25营收(亿)} & \textbf{营收YoY} & \textbf{26Q1净利(亿)} & \textbf{净利YoY} & \textbf{毛利率Q1} & \textbf{存货(亿)} & \textbf{经营CF(亿)} & \textbf{ROE} & \textbf{核心看点} \\
\midrule
\rowcolor{riskgreen!8}
\textbf{江波龙} & 227.66 & +30.4\% & \textbf{38.62} & \textbf{+2644\%} & \textbf{55.53\%} & \textbf{116.78} & -12.01 & 18.1\% & \textbf{\astockstar{} 存储涨价弹性最大} \\
\rowcolor{riskgreen!4}
北京君正 & 47.41 & +12.5\% & 3.19 & +332\% & 43.49\% & 29.75 & 6.23 & 3.06\% & 车载存储+涨价确认 \\
\rowcolor{riskgreen!4}
兆易创新 & 92.03 & +25.1\% & 14.61 & +523\% & 57.08\% & 34.01 & 21.29 & 9.30\% & NOR Flash+利基DRAM涨价 \\
\rowcolor{riskgreen!4}
澜起科技 & 54.56 & +49.9\% & 8.47 & +61.3\% & 69.80\% & 8.30 & 20.22 & 18.25\% & DDR5 RCD全球第一+CXL \\
北方华创 & 393.53 & +30.9\% & 16.35 & +3.4\% & 40.77\% & 286.27 & 21.33 & 16.41\% & 设备平台龙头，订单饱满 \\
中微公司 & 123.85 & +36.6\% & 9.30 & +197\% & 39.89\% & 20.64 & 22.80 & 9.97\% & 刻蚀设备龙头+海外突破 \\
拓荆科技 & 65.19 & +58.9\% & 5.71 & +488\% & 41.69\% & $\sim$10 & 35.90 & 15.77\% & ALD/CVD最快增长，\textbf{停牌} \\
沪硅产业 & 37.16 & +9.7\% & -4.83 & -132\% & -11.85\% & 19.66 & -5.59 & -12.68\% & 大硅片国产替代，\textbf{PE屏蔽} \\
深科技 & 157.47 & +6.2\% & 2.42 & +35\% & 17.07\% & 28.67 & 17.65 & 9.09\% & 封测+模组双主业 \\
长电科技 & 388.71 & +8.1\% & 2.90 & +43\% & 14.55\% & 44.08 & 46.52 & 5.46\% & 先进封测全球第三 \\
通富微电 & 279.21 & +16.9\% & 3.29 & +225\% & 13.32\% & 48.11 & 69.66 & 8.08\% & AMD绑定+定增获批 \\
南大光电 & 25.85 & +9.9\% & 1.24 & +30\% & 41.86\% & 8.11 & 7.33 & 9.19\% & MO源龙头+前驱体 \\
\midrule
\rowcolor{panelgrey}
\textbf{合计/均值} & \textbf{1,893} & -- & \textbf{91.2} & -- & -- & \textbf{644} & -- & -- & -- \\
\bottomrule
\end{tabularx}
\end{exhibitbox}
\sourcenote{数据来自各公司2025年报+2026Q1季报（2026-04至05披露）。江波龙中报预告来自2026-07-03公告。存货口径：江波龙/北京君正/沪硅为2025年末；其余为2026Q1。\textbf{关键洞察}：存储涨价周期中模组厂(江波龙)利润率弹性最大（毛利率19\%$\to$56\%），设备厂(拓荆/中微)营收增速最快，材料厂(沪硅)仍在亏损爬坡期。}

\vspace{0.4em}
\noindent\textbf{跨公司比较的投资含义}：

\begin{enumerate}[label=\protect\gmark]
  \item \textbf{库存/营收比}：江波龙0.51x（2025年末备货最积极$\to$涨价弹性最大）、拓荆约0.15x（设备交付周期长，在手订单饱满度以预收款体现）、长电0.48x（封测周转快）
  \item \textbf{毛利率变化（Q1 vs 2025FY）}：江波龙+36.1pp（最显著）、兆易+16.9pp、北京君正+9.4pp、澜起+7.6pp——存储涨价直接受益者
  \item \textbf{现金流质量}：通富微电经营CF/净利=5.7x（封测先收款后结算模式）、拓荆3.87x（设备预收款）、南大光电2.29x（前驱体盈利质量高）；模组厂(江波龙)备货期现金流为负但利润弹性大
  \item \textbf{重资产vs轻资产}：沪硅(141亿)/长电(约180亿)/通富(约120亿)/北方华创(76亿)重资产（CAPEX大，产能壁垒高）；澜起/兆易/君正轻资产（Fabless，ROE弹性大）
  \item \textbf{AI存储敞口强度}：极高（江波龙、澜起）$>$高（北方华创、中微、拓荆）$>$中高（北京君正）$>$中（沪硅、长电、兆易）$>$中低（深科技、通富、南大光电）
\end{enumerate}